% ============================================================================
% preamble.tex -- stil, yazı tipleri, paketler, renkler, kutu/tablo stilleri.
% main.tex tarafından % ============================================================================
% preamble.tex -- stil, yazı tipleri, paketler, renkler, kutu/tablo stilleri.
% main.tex tarafından % ============================================================================
% preamble.tex -- stil, yazı tipleri, paketler, renkler, kutu/tablo stilleri.
% main.tex tarafından % ============================================================================
% preamble.tex -- stil, yazı tipleri, paketler, renkler, kutu/tablo stilleri.
% main.tex tarafından \input{preamble} ile çağrılır.
% ============================================================================

% ---- Sayfa düzeni ----------------------------------------------------------
\KOMAoptions{
  fontsize=11pt,
  paper=a4,
  DIV=12,
  parskip=half,
  headsepline=true,
  footsepline=false,
  numbers=noenddot,
  chapterprefix=false,
}

% ---- Türkçe + yazı tipleri (XeLaTeX) ---------------------------------------
\usepackage{fontspec}
\usepackage{polyglossia}
\setmainlanguage{turkish}

% TeX Gyre ailesi: URW/Adobe klasik fontlarının serbestçe dağıtılabilir
% OpenType karşılıkları. Herhangi bir TeX Live kurulumunda (texlive-fonts-extra
% ya da tek başına fonts-texgyre paketiyle) bulunur; internetten ayrı bir
% font indirmeye gerek yoktur.
\setmainfont{TeX Gyre Pagella}                      % gövde metni (Palatino benzeri)
\setsansfont{TeX Gyre Heros}                        % başlıklar (Helvetica benzeri)
\setmonofont{TeX Gyre Cursor}[Scale=0.92]           % kod/komut satırları (Courier benzeri)

% ---- Renkler ----------------------------------------------------------------
% Deneyler.tex ve EXPERIMENT_LOG.md'deki grafiklerle aynı paletten:
% koyu lacivert ana renk, mercan ikincil vurgu, nötr griler.
\usepackage{xcolor}
\definecolor{anarenk}{HTML}{1B3A6B}       % başlıklar, çizgiler
\definecolor{vurgu}{HTML}{C8551A}         % dikkat çekilecek sayılar
\definecolor{yesilvurgu}{HTML}{1E7A52}    % "ölçüldü / geçti" göstergeleri
\definecolor{grimetin}{HTML}{4A4A46}
\definecolor{acikgri}{HTML}{EDEBE4}
\definecolor{kutucerceve}{HTML}{B9B6A9}

% ---- Başlık stilleri ---------------------------------------------------------
% (titlesec kasıtlı olarak kullanılmıyor: KOMA-Script kendi
% \addtokomafont mekanizmasıyla stil veriyor, titlesec ile birlikte
% yüklenmesi scrlayer-scrpage'le çakışıp uyarı üretiyordu.)
\addtokomafont{chapter}{\color{anarenk}\sffamily}
\addtokomafont{section}{\color{anarenk}\sffamily}
\addtokomafont{subsection}{\color{anarenk}\sffamily}
\addtokomafont{subsubsection}{\sffamily\bfseries}
\addtokomafont{disposition}{\rmfamily}     % KOMA varsayılanını geçersiz kılıp yukarıdakileri kullan
\addtokomafont{pagehead}{\sffamily\small\color{grimetin}}
\addtokomafont{pagefoot}{\sffamily\small\color{grimetin}}

% ---- Tablolar, şekiller, kod -------------------------------------------------
\usepackage{booktabs}
\usepackage{array}
\usepackage{longtable}
\usepackage{graphicx}
\usepackage{caption}
\captionsetup{font=small, labelfont={sf,bf,color=anarenk}, margin=1em}
\usepackage{subcaption}
\usepackage{etoolbox}
\usepackage{float}
\usepackage{tikz}
\usetikzlibrary{arrows.meta, positioning, fit, backgrounds, calc}
\usepackage{listings}
\lstset{
  basicstyle=\ttfamily\small,
  breaklines=true,
  columns=fullflexible,
  frame=single,
  rulecolor=\color{kutucerceve},
  backgroundcolor=\color{acikgri!45},
  captionpos=b,
}

% ---- Vurgu/uyarı kutuları ----------------------------------------------------
\usepackage[most]{tcolorbox}
\tcbuselibrary{skins,breakable}

% Ölçülmüş, alıntılanabilir sonuç kutusu.
\newtcolorbox{olcum}[1][]{
  enhanced, breakable, colback=acikgri!35, colframe=anarenk,
  boxrule=0.6pt, left=2mm, right=2mm, top=1mm, bottom=1mm,
  fonttitle=\sffamily\bfseries, coltitle=white,
  colbacktitle=anarenk, title={Ölçüldü (cihazda)}, #1,
}

% CPU'da / yapısal olarak doğrulanmış, ama Orin'de henüz alınmamış.
\newtcolorbox{yapisal}[1][]{
  enhanced, breakable, colback=white, colframe=grimetin,
  boxrule=0.5pt, left=2mm, right=2mm, top=1mm, bottom=1mm,
  fonttitle=\sffamily\bfseries, coltitle=white,
  colbacktitle=grimetin, title={CPU'da, yapısal}, #1,
}

% Henüz ölçülmemiş -- rapor tamamlanmadan doldurulması gereken yer.
\newtcolorbox{eksik}[1][]{
  enhanced, breakable, colback=vurgu!8, colframe=vurgu,
  boxrule=0.6pt, left=2mm, right=2mm, top=1mm, bottom=1mm,
  fonttitle=\sffamily\bfseries, coltitle=white,
  colbacktitle=vurgu, title={TODO -- cihazda ölçülüp doldurulacak}, #1,
}

% Görsel/figür yer tutucusu: kullanıcı kendi görselini yerel olarak
% ekleyecek, bu kutu onun yerini işaretler ve derlemeyi hiçbir zaman
% bozmaz (eksik bir \includegraphics dosyası derleme hatası verirdi).
\newcommand{\gorselyertutucu}[2][0.85\linewidth]{%
  \begin{center}
  \fbox{\begin{minipage}{#1}
    \centering\vspace{2.2cm}
    {\sffamily\bfseries\color{grimetin} GÖRSEL BURAYA}\\[2mm]
    {\small\color{grimetin} #2}
    \vspace{2.2cm}
  \end{minipage}}
  \end{center}
}

% Dört görseli 2x2 ızgarada, her birinin altında kendi kısa alt yazısıyla
% gösteren şablon: crop/full karşılaştırması, dörtlü grafik sonuçları vb.
% için tekrar kullanılabilir. Dosya yolu boş bırakılan hücreye
% \gorselyertutucu çizilir, derleme hiçbir zaman bozulmaz.
\newcommand{\gorselhucre}[3][0.47\linewidth]{%
  % #1 (ops.): hücre genişliği; #2: dosya yolu (boşsa yer tutucu); #3: alt yazı
  \begin{subfigure}[t]{#1}
    \centering
    \ifstrempty{#2}{\gorselyertutucu[0.92\linewidth]{#3}}{%
      \includegraphics[width=\linewidth]{#2}}
    \caption{#3}
  \end{subfigure}%
}
% #1--#4: dosya yolları (soldan sağa, üstten alta; boşsa yer tutucu);
% #5--#8: alt yazılar. \caption ve \label çağıran tarafta kalır -- bunu
% her zaman \begin{figure}[H] ... \end{figure} içine sarın (bkz. kullanım
% örnekleri deneyler.tex içinde \S{Gerçek Kayıtlar Üzerinde Girdi
% Konfigürasyonu}'nda).
\newcommand{\dortlugorsel}[8]{%
  \gorselhucre{#1}{#5}\hfill\gorselhucre{#2}{#6}\\[2mm]
  \gorselhucre{#3}{#7}\hfill\gorselhucre{#4}{#8}%
}

% ---- Bağlantılar, referanslar ------------------------------------------------
\usepackage{hyperref}
\hypersetup{
  colorlinks=true,
  linkcolor=anarenk,
  citecolor=yesilvurgu,
  urlcolor=vurgu,
  pdftitle={EdgeTAM'ın Jetson AGX Orin Üzerinde TensorRT ile Hızlandırılması},
  pdfauthor={Yiğit Kaya Bağcı},
}
% english: cleveref'in kendi dahili Türkçe ad tablosu polyglossia'nın
% Türkçe büyük/küçük harf kurallarıyla (noktasız ı / noktalı İ) çakışıp
% \begin{document}'te "Undefined control sequence" hatası veriyordu --
% bilinen bir cleveref+polyglossia-turkish etkileşim sorunu. english
% seçeneği cleveref'in dahili tablosunu devre dışı bırakır; görünen adlar
% zaten aşağıdaki \crefname tanımlarından geliyor.
\usepackage[capitalize, nameinlink, english]{cleveref}
\crefname{table}{Tablo}{Tablolar}
\Crefname{table}{Tablo}{Tablolar}
\crefname{figure}{Şekil}{Şekiller}
\Crefname{figure}{Şekil}{Şekiller}
\crefname{chapter}{Bölüm}{Bölümler}
\Crefname{chapter}{Bölüm}{Bölümler}
\crefname{section}{Bölüm}{Bölümler}
\Crefname{section}{Bölüm}{Bölümler}

% ---- Kaynakça -----------------------------------------------------------------
\usepackage[
  backend=biber,
  style=numeric,
  sorting=none,
  giveninits=true,
]{biblatex}
\addbibresource{kaynaklar.bib}

% ---- Matematik ------------------------------------------------------------
\usepackage{amsmath}
\usepackage{amssymb}

% \theauthor / \thedate: \title, \author, \date içeriğine \maketitle
% dışında da erişim (özel başlık sayfası için).
\usepackage{titling}

% ---- Üstbilgi/altbilgi ------------------------------------------------------
\usepackage[automark]{scrlayer-scrpage}
\pagestyle{scrheadings}
\clearpairofpagestyles
\ihead{\headmark}
\ohead{\thepage}
\ifoot{}
\ofoot{EdgeTAM / TensorRT -- Jetson AGX Orin}
 ile çağrılır.
% ============================================================================

% ---- Sayfa düzeni ----------------------------------------------------------
\KOMAoptions{
  fontsize=11pt,
  paper=a4,
  DIV=12,
  parskip=half,
  headsepline=true,
  footsepline=false,
  numbers=noenddot,
  chapterprefix=false,
}

% ---- Türkçe + yazı tipleri (XeLaTeX) ---------------------------------------
\usepackage{fontspec}
\usepackage{polyglossia}
\setmainlanguage{turkish}

% TeX Gyre ailesi: URW/Adobe klasik fontlarının serbestçe dağıtılabilir
% OpenType karşılıkları. Herhangi bir TeX Live kurulumunda (texlive-fonts-extra
% ya da tek başına fonts-texgyre paketiyle) bulunur; internetten ayrı bir
% font indirmeye gerek yoktur.
\setmainfont{TeX Gyre Pagella}                      % gövde metni (Palatino benzeri)
\setsansfont{TeX Gyre Heros}                        % başlıklar (Helvetica benzeri)
\setmonofont{TeX Gyre Cursor}[Scale=0.92]           % kod/komut satırları (Courier benzeri)

% ---- Renkler ----------------------------------------------------------------
% Deneyler.tex ve EXPERIMENT_LOG.md'deki grafiklerle aynı paletten:
% koyu lacivert ana renk, mercan ikincil vurgu, nötr griler.
\usepackage{xcolor}
\definecolor{anarenk}{HTML}{1B3A6B}       % başlıklar, çizgiler
\definecolor{vurgu}{HTML}{C8551A}         % dikkat çekilecek sayılar
\definecolor{yesilvurgu}{HTML}{1E7A52}    % "ölçüldü / geçti" göstergeleri
\definecolor{grimetin}{HTML}{4A4A46}
\definecolor{acikgri}{HTML}{EDEBE4}
\definecolor{kutucerceve}{HTML}{B9B6A9}

% ---- Başlık stilleri ---------------------------------------------------------
% (titlesec kasıtlı olarak kullanılmıyor: KOMA-Script kendi
% \addtokomafont mekanizmasıyla stil veriyor, titlesec ile birlikte
% yüklenmesi scrlayer-scrpage'le çakışıp uyarı üretiyordu.)
\addtokomafont{chapter}{\color{anarenk}\sffamily}
\addtokomafont{section}{\color{anarenk}\sffamily}
\addtokomafont{subsection}{\color{anarenk}\sffamily}
\addtokomafont{subsubsection}{\sffamily\bfseries}
\addtokomafont{disposition}{\rmfamily}     % KOMA varsayılanını geçersiz kılıp yukarıdakileri kullan
\addtokomafont{pagehead}{\sffamily\small\color{grimetin}}
\addtokomafont{pagefoot}{\sffamily\small\color{grimetin}}

% ---- Tablolar, şekiller, kod -------------------------------------------------
\usepackage{booktabs}
\usepackage{array}
\usepackage{longtable}
\usepackage{graphicx}
\usepackage{caption}
\captionsetup{font=small, labelfont={sf,bf,color=anarenk}, margin=1em}
\usepackage{subcaption}
\usepackage{etoolbox}
\usepackage{float}
\usepackage{tikz}
\usetikzlibrary{arrows.meta, positioning, fit, backgrounds, calc}
\usepackage{listings}
\lstset{
  basicstyle=\ttfamily\small,
  breaklines=true,
  columns=fullflexible,
  frame=single,
  rulecolor=\color{kutucerceve},
  backgroundcolor=\color{acikgri!45},
  captionpos=b,
}

% ---- Vurgu/uyarı kutuları ----------------------------------------------------
\usepackage[most]{tcolorbox}
\tcbuselibrary{skins,breakable}

% Ölçülmüş, alıntılanabilir sonuç kutusu.
\newtcolorbox{olcum}[1][]{
  enhanced, breakable, colback=acikgri!35, colframe=anarenk,
  boxrule=0.6pt, left=2mm, right=2mm, top=1mm, bottom=1mm,
  fonttitle=\sffamily\bfseries, coltitle=white,
  colbacktitle=anarenk, title={Ölçüldü (cihazda)}, #1,
}

% CPU'da / yapısal olarak doğrulanmış, ama Orin'de henüz alınmamış.
\newtcolorbox{yapisal}[1][]{
  enhanced, breakable, colback=white, colframe=grimetin,
  boxrule=0.5pt, left=2mm, right=2mm, top=1mm, bottom=1mm,
  fonttitle=\sffamily\bfseries, coltitle=white,
  colbacktitle=grimetin, title={CPU'da, yapısal}, #1,
}

% Henüz ölçülmemiş -- rapor tamamlanmadan doldurulması gereken yer.
\newtcolorbox{eksik}[1][]{
  enhanced, breakable, colback=vurgu!8, colframe=vurgu,
  boxrule=0.6pt, left=2mm, right=2mm, top=1mm, bottom=1mm,
  fonttitle=\sffamily\bfseries, coltitle=white,
  colbacktitle=vurgu, title={TODO -- cihazda ölçülüp doldurulacak}, #1,
}

% Görsel/figür yer tutucusu: kullanıcı kendi görselini yerel olarak
% ekleyecek, bu kutu onun yerini işaretler ve derlemeyi hiçbir zaman
% bozmaz (eksik bir \includegraphics dosyası derleme hatası verirdi).
\newcommand{\gorselyertutucu}[2][0.85\linewidth]{%
  \begin{center}
  \fbox{\begin{minipage}{#1}
    \centering\vspace{2.2cm}
    {\sffamily\bfseries\color{grimetin} GÖRSEL BURAYA}\\[2mm]
    {\small\color{grimetin} #2}
    \vspace{2.2cm}
  \end{minipage}}
  \end{center}
}

% Dört görseli 2x2 ızgarada, her birinin altında kendi kısa alt yazısıyla
% gösteren şablon: crop/full karşılaştırması, dörtlü grafik sonuçları vb.
% için tekrar kullanılabilir. Dosya yolu boş bırakılan hücreye
% \gorselyertutucu çizilir, derleme hiçbir zaman bozulmaz.
\newcommand{\gorselhucre}[3][0.47\linewidth]{%
  % #1 (ops.): hücre genişliği; #2: dosya yolu (boşsa yer tutucu); #3: alt yazı
  \begin{subfigure}[t]{#1}
    \centering
    \ifstrempty{#2}{\gorselyertutucu[0.92\linewidth]{#3}}{%
      \includegraphics[width=\linewidth]{#2}}
    \caption{#3}
  \end{subfigure}%
}
% #1--#4: dosya yolları (soldan sağa, üstten alta; boşsa yer tutucu);
% #5--#8: alt yazılar. \caption ve \label çağıran tarafta kalır -- bunu
% her zaman \begin{figure}[H] ... \end{figure} içine sarın (bkz. kullanım
% örnekleri deneyler.tex içinde \S{Gerçek Kayıtlar Üzerinde Girdi
% Konfigürasyonu}'nda).
\newcommand{\dortlugorsel}[8]{%
  \gorselhucre{#1}{#5}\hfill\gorselhucre{#2}{#6}\\[2mm]
  \gorselhucre{#3}{#7}\hfill\gorselhucre{#4}{#8}%
}

% ---- Bağlantılar, referanslar ------------------------------------------------
\usepackage{hyperref}
\hypersetup{
  colorlinks=true,
  linkcolor=anarenk,
  citecolor=yesilvurgu,
  urlcolor=vurgu,
  pdftitle={EdgeTAM'ın Jetson AGX Orin Üzerinde TensorRT ile Hızlandırılması},
  pdfauthor={Yiğit Kaya Bağcı},
}
% english: cleveref'in kendi dahili Türkçe ad tablosu polyglossia'nın
% Türkçe büyük/küçük harf kurallarıyla (noktasız ı / noktalı İ) çakışıp
% \begin{document}'te "Undefined control sequence" hatası veriyordu --
% bilinen bir cleveref+polyglossia-turkish etkileşim sorunu. english
% seçeneği cleveref'in dahili tablosunu devre dışı bırakır; görünen adlar
% zaten aşağıdaki \crefname tanımlarından geliyor.
\usepackage[capitalize, nameinlink, english]{cleveref}
\crefname{table}{Tablo}{Tablolar}
\Crefname{table}{Tablo}{Tablolar}
\crefname{figure}{Şekil}{Şekiller}
\Crefname{figure}{Şekil}{Şekiller}
\crefname{chapter}{Bölüm}{Bölümler}
\Crefname{chapter}{Bölüm}{Bölümler}
\crefname{section}{Bölüm}{Bölümler}
\Crefname{section}{Bölüm}{Bölümler}

% ---- Kaynakça -----------------------------------------------------------------
\usepackage[
  backend=biber,
  style=numeric,
  sorting=none,
  giveninits=true,
]{biblatex}
\addbibresource{kaynaklar.bib}

% ---- Matematik ------------------------------------------------------------
\usepackage{amsmath}
\usepackage{amssymb}

% \theauthor / \thedate: \title, \author, \date içeriğine \maketitle
% dışında da erişim (özel başlık sayfası için).
\usepackage{titling}

% ---- Üstbilgi/altbilgi ------------------------------------------------------
\usepackage[automark]{scrlayer-scrpage}
\pagestyle{scrheadings}
\clearpairofpagestyles
\ihead{\headmark}
\ohead{\thepage}
\ifoot{}
\ofoot{EdgeTAM / TensorRT -- Jetson AGX Orin}
 ile çağrılır.
% ============================================================================

% ---- Sayfa düzeni ----------------------------------------------------------
\KOMAoptions{
  fontsize=11pt,
  paper=a4,
  DIV=12,
  parskip=half,
  headsepline=true,
  footsepline=false,
  numbers=noenddot,
  chapterprefix=false,
}

% ---- Türkçe + yazı tipleri (XeLaTeX) ---------------------------------------
\usepackage{fontspec}
\usepackage{polyglossia}
\setmainlanguage{turkish}

% TeX Gyre ailesi: URW/Adobe klasik fontlarının serbestçe dağıtılabilir
% OpenType karşılıkları. Herhangi bir TeX Live kurulumunda (texlive-fonts-extra
% ya da tek başına fonts-texgyre paketiyle) bulunur; internetten ayrı bir
% font indirmeye gerek yoktur.
\setmainfont{TeX Gyre Pagella}                      % gövde metni (Palatino benzeri)
\setsansfont{TeX Gyre Heros}                        % başlıklar (Helvetica benzeri)
\setmonofont{TeX Gyre Cursor}[Scale=0.92]           % kod/komut satırları (Courier benzeri)

% ---- Renkler ----------------------------------------------------------------
% Deneyler.tex ve EXPERIMENT_LOG.md'deki grafiklerle aynı paletten:
% koyu lacivert ana renk, mercan ikincil vurgu, nötr griler.
\usepackage{xcolor}
\definecolor{anarenk}{HTML}{1B3A6B}       % başlıklar, çizgiler
\definecolor{vurgu}{HTML}{C8551A}         % dikkat çekilecek sayılar
\definecolor{yesilvurgu}{HTML}{1E7A52}    % "ölçüldü / geçti" göstergeleri
\definecolor{grimetin}{HTML}{4A4A46}
\definecolor{acikgri}{HTML}{EDEBE4}
\definecolor{kutucerceve}{HTML}{B9B6A9}

% ---- Başlık stilleri ---------------------------------------------------------
% (titlesec kasıtlı olarak kullanılmıyor: KOMA-Script kendi
% \addtokomafont mekanizmasıyla stil veriyor, titlesec ile birlikte
% yüklenmesi scrlayer-scrpage'le çakışıp uyarı üretiyordu.)
\addtokomafont{chapter}{\color{anarenk}\sffamily}
\addtokomafont{section}{\color{anarenk}\sffamily}
\addtokomafont{subsection}{\color{anarenk}\sffamily}
\addtokomafont{subsubsection}{\sffamily\bfseries}
\addtokomafont{disposition}{\rmfamily}     % KOMA varsayılanını geçersiz kılıp yukarıdakileri kullan
\addtokomafont{pagehead}{\sffamily\small\color{grimetin}}
\addtokomafont{pagefoot}{\sffamily\small\color{grimetin}}

% ---- Tablolar, şekiller, kod -------------------------------------------------
\usepackage{booktabs}
\usepackage{array}
\usepackage{longtable}
\usepackage{graphicx}
\usepackage{caption}
\captionsetup{font=small, labelfont={sf,bf,color=anarenk}, margin=1em}
\usepackage{subcaption}
\usepackage{etoolbox}
\usepackage{float}
\usepackage{tikz}
\usetikzlibrary{arrows.meta, positioning, fit, backgrounds, calc}
\usepackage{listings}
\lstset{
  basicstyle=\ttfamily\small,
  breaklines=true,
  columns=fullflexible,
  frame=single,
  rulecolor=\color{kutucerceve},
  backgroundcolor=\color{acikgri!45},
  captionpos=b,
}

% ---- Vurgu/uyarı kutuları ----------------------------------------------------
\usepackage[most]{tcolorbox}
\tcbuselibrary{skins,breakable}

% Ölçülmüş, alıntılanabilir sonuç kutusu.
\newtcolorbox{olcum}[1][]{
  enhanced, breakable, colback=acikgri!35, colframe=anarenk,
  boxrule=0.6pt, left=2mm, right=2mm, top=1mm, bottom=1mm,
  fonttitle=\sffamily\bfseries, coltitle=white,
  colbacktitle=anarenk, title={Ölçüldü (cihazda)}, #1,
}

% CPU'da / yapısal olarak doğrulanmış, ama Orin'de henüz alınmamış.
\newtcolorbox{yapisal}[1][]{
  enhanced, breakable, colback=white, colframe=grimetin,
  boxrule=0.5pt, left=2mm, right=2mm, top=1mm, bottom=1mm,
  fonttitle=\sffamily\bfseries, coltitle=white,
  colbacktitle=grimetin, title={CPU'da, yapısal}, #1,
}

% Henüz ölçülmemiş -- rapor tamamlanmadan doldurulması gereken yer.
\newtcolorbox{eksik}[1][]{
  enhanced, breakable, colback=vurgu!8, colframe=vurgu,
  boxrule=0.6pt, left=2mm, right=2mm, top=1mm, bottom=1mm,
  fonttitle=\sffamily\bfseries, coltitle=white,
  colbacktitle=vurgu, title={TODO -- cihazda ölçülüp doldurulacak}, #1,
}

% Görsel/figür yer tutucusu: kullanıcı kendi görselini yerel olarak
% ekleyecek, bu kutu onun yerini işaretler ve derlemeyi hiçbir zaman
% bozmaz (eksik bir \includegraphics dosyası derleme hatası verirdi).
\newcommand{\gorselyertutucu}[2][0.85\linewidth]{%
  \begin{center}
  \fbox{\begin{minipage}{#1}
    \centering\vspace{2.2cm}
    {\sffamily\bfseries\color{grimetin} GÖRSEL BURAYA}\\[2mm]
    {\small\color{grimetin} #2}
    \vspace{2.2cm}
  \end{minipage}}
  \end{center}
}

% Dört görseli 2x2 ızgarada, her birinin altında kendi kısa alt yazısıyla
% gösteren şablon: crop/full karşılaştırması, dörtlü grafik sonuçları vb.
% için tekrar kullanılabilir. Dosya yolu boş bırakılan hücreye
% \gorselyertutucu çizilir, derleme hiçbir zaman bozulmaz.
\newcommand{\gorselhucre}[3][0.47\linewidth]{%
  % #1 (ops.): hücre genişliği; #2: dosya yolu (boşsa yer tutucu); #3: alt yazı
  \begin{subfigure}[t]{#1}
    \centering
    \ifstrempty{#2}{\gorselyertutucu[0.92\linewidth]{#3}}{%
      \includegraphics[width=\linewidth]{#2}}
    \caption{#3}
  \end{subfigure}%
}
% #1--#4: dosya yolları (soldan sağa, üstten alta; boşsa yer tutucu);
% #5--#8: alt yazılar. \caption ve \label çağıran tarafta kalır -- bunu
% her zaman \begin{figure}[H] ... \end{figure} içine sarın (bkz. kullanım
% örnekleri deneyler.tex içinde \S{Gerçek Kayıtlar Üzerinde Girdi
% Konfigürasyonu}'nda).
\newcommand{\dortlugorsel}[8]{%
  \gorselhucre{#1}{#5}\hfill\gorselhucre{#2}{#6}\\[2mm]
  \gorselhucre{#3}{#7}\hfill\gorselhucre{#4}{#8}%
}

% ---- Bağlantılar, referanslar ------------------------------------------------
\usepackage{hyperref}
\hypersetup{
  colorlinks=true,
  linkcolor=anarenk,
  citecolor=yesilvurgu,
  urlcolor=vurgu,
  pdftitle={EdgeTAM'ın Jetson AGX Orin Üzerinde TensorRT ile Hızlandırılması},
  pdfauthor={Yiğit Kaya Bağcı},
}
% english: cleveref'in kendi dahili Türkçe ad tablosu polyglossia'nın
% Türkçe büyük/küçük harf kurallarıyla (noktasız ı / noktalı İ) çakışıp
% \begin{document}'te "Undefined control sequence" hatası veriyordu --
% bilinen bir cleveref+polyglossia-turkish etkileşim sorunu. english
% seçeneği cleveref'in dahili tablosunu devre dışı bırakır; görünen adlar
% zaten aşağıdaki \crefname tanımlarından geliyor.
\usepackage[capitalize, nameinlink, english]{cleveref}
\crefname{table}{Tablo}{Tablolar}
\Crefname{table}{Tablo}{Tablolar}
\crefname{figure}{Şekil}{Şekiller}
\Crefname{figure}{Şekil}{Şekiller}
\crefname{chapter}{Bölüm}{Bölümler}
\Crefname{chapter}{Bölüm}{Bölümler}
\crefname{section}{Bölüm}{Bölümler}
\Crefname{section}{Bölüm}{Bölümler}

% ---- Kaynakça -----------------------------------------------------------------
\usepackage[
  backend=biber,
  style=numeric,
  sorting=none,
  giveninits=true,
]{biblatex}
\addbibresource{kaynaklar.bib}

% ---- Matematik ------------------------------------------------------------
\usepackage{amsmath}
\usepackage{amssymb}

% \theauthor / \thedate: \title, \author, \date içeriğine \maketitle
% dışında da erişim (özel başlık sayfası için).
\usepackage{titling}

% ---- Üstbilgi/altbilgi ------------------------------------------------------
\usepackage[automark]{scrlayer-scrpage}
\pagestyle{scrheadings}
\clearpairofpagestyles
\ihead{\headmark}
\ohead{\thepage}
\ifoot{}
\ofoot{EdgeTAM / TensorRT -- Jetson AGX Orin}
 ile çağrılır.
% ============================================================================

% ---- Sayfa düzeni ----------------------------------------------------------
\KOMAoptions{
  fontsize=11pt,
  paper=a4,
  DIV=12,
  parskip=half,
  headsepline=true,
  footsepline=false,
  numbers=noenddot,
  chapterprefix=false,
}

% ---- Türkçe + yazı tipleri (XeLaTeX) ---------------------------------------
\usepackage{fontspec}
\usepackage{polyglossia}
\setmainlanguage{turkish}

% TeX Gyre ailesi: URW/Adobe klasik fontlarının serbestçe dağıtılabilir
% OpenType karşılıkları. Herhangi bir TeX Live kurulumunda (texlive-fonts-extra
% ya da tek başına fonts-texgyre paketiyle) bulunur; internetten ayrı bir
% font indirmeye gerek yoktur.
\setmainfont{TeX Gyre Pagella}                      % gövde metni (Palatino benzeri)
\setsansfont{TeX Gyre Heros}                        % başlıklar (Helvetica benzeri)
\setmonofont{TeX Gyre Cursor}[Scale=0.92]           % kod/komut satırları (Courier benzeri)

% ---- Renkler ----------------------------------------------------------------
% Deneyler.tex ve EXPERIMENT_LOG.md'deki grafiklerle aynı paletten:
% koyu lacivert ana renk, mercan ikincil vurgu, nötr griler.
\usepackage{xcolor}
\definecolor{anarenk}{HTML}{1B3A6B}       % başlıklar, çizgiler
\definecolor{vurgu}{HTML}{C8551A}         % dikkat çekilecek sayılar
\definecolor{yesilvurgu}{HTML}{1E7A52}    % "ölçüldü / geçti" göstergeleri
\definecolor{grimetin}{HTML}{4A4A46}
\definecolor{acikgri}{HTML}{EDEBE4}
\definecolor{kutucerceve}{HTML}{B9B6A9}

% ---- Başlık stilleri ---------------------------------------------------------
% (titlesec kasıtlı olarak kullanılmıyor: KOMA-Script kendi
% \addtokomafont mekanizmasıyla stil veriyor, titlesec ile birlikte
% yüklenmesi scrlayer-scrpage'le çakışıp uyarı üretiyordu.)
\addtokomafont{chapter}{\color{anarenk}\sffamily}
\addtokomafont{section}{\color{anarenk}\sffamily}
\addtokomafont{subsection}{\color{anarenk}\sffamily}
\addtokomafont{subsubsection}{\sffamily\bfseries}
\addtokomafont{disposition}{\rmfamily}     % KOMA varsayılanını geçersiz kılıp yukarıdakileri kullan
\addtokomafont{pagehead}{\sffamily\small\color{grimetin}}
\addtokomafont{pagefoot}{\sffamily\small\color{grimetin}}

% ---- Tablolar, şekiller, kod -------------------------------------------------
\usepackage{booktabs}
\usepackage{array}
\usepackage{longtable}
\usepackage{graphicx}
\usepackage{caption}
\captionsetup{font=small, labelfont={sf,bf,color=anarenk}, margin=1em}
\usepackage{subcaption}
\usepackage{etoolbox}
\usepackage{float}
\usepackage{tikz}
\usetikzlibrary{arrows.meta, positioning, fit, backgrounds, calc}
\usepackage{listings}
\lstset{
  basicstyle=\ttfamily\small,
  breaklines=true,
  columns=fullflexible,
  frame=single,
  rulecolor=\color{kutucerceve},
  backgroundcolor=\color{acikgri!45},
  captionpos=b,
}

% ---- Vurgu/uyarı kutuları ----------------------------------------------------
\usepackage[most]{tcolorbox}
\tcbuselibrary{skins,breakable}

% Ölçülmüş, alıntılanabilir sonuç kutusu.
\newtcolorbox{olcum}[1][]{
  enhanced, breakable, colback=acikgri!35, colframe=anarenk,
  boxrule=0.6pt, left=2mm, right=2mm, top=1mm, bottom=1mm,
  fonttitle=\sffamily\bfseries, coltitle=white,
  colbacktitle=anarenk, title={Ölçüldü (cihazda)}, #1,
}

% CPU'da / yapısal olarak doğrulanmış, ama Orin'de henüz alınmamış.
\newtcolorbox{yapisal}[1][]{
  enhanced, breakable, colback=white, colframe=grimetin,
  boxrule=0.5pt, left=2mm, right=2mm, top=1mm, bottom=1mm,
  fonttitle=\sffamily\bfseries, coltitle=white,
  colbacktitle=grimetin, title={CPU'da, yapısal}, #1,
}

% Henüz ölçülmemiş -- rapor tamamlanmadan doldurulması gereken yer.
\newtcolorbox{eksik}[1][]{
  enhanced, breakable, colback=vurgu!8, colframe=vurgu,
  boxrule=0.6pt, left=2mm, right=2mm, top=1mm, bottom=1mm,
  fonttitle=\sffamily\bfseries, coltitle=white,
  colbacktitle=vurgu, title={TODO -- cihazda ölçülüp doldurulacak}, #1,
}

% Görsel/figür yer tutucusu: kullanıcı kendi görselini yerel olarak
% ekleyecek, bu kutu onun yerini işaretler ve derlemeyi hiçbir zaman
% bozmaz (eksik bir \includegraphics dosyası derleme hatası verirdi).
\newcommand{\gorselyertutucu}[2][0.85\linewidth]{%
  \begin{center}
  \fbox{\begin{minipage}{#1}
    \centering\vspace{2.2cm}
    {\sffamily\bfseries\color{grimetin} GÖRSEL BURAYA}\\[2mm]
    {\small\color{grimetin} #2}
    \vspace{2.2cm}
  \end{minipage}}
  \end{center}
}

% Dört görseli 2x2 ızgarada, her birinin altında kendi kısa alt yazısıyla
% gösteren şablon: crop/full karşılaştırması, dörtlü grafik sonuçları vb.
% için tekrar kullanılabilir. Dosya yolu boş bırakılan hücreye
% \gorselyertutucu çizilir, derleme hiçbir zaman bozulmaz.
\newcommand{\gorselhucre}[3][0.47\linewidth]{%
  % #1 (ops.): hücre genişliği; #2: dosya yolu (boşsa yer tutucu); #3: alt yazı
  \begin{subfigure}[t]{#1}
    \centering
    \ifstrempty{#2}{\gorselyertutucu[0.92\linewidth]{#3}}{%
      \includegraphics[width=\linewidth]{#2}}
    \caption{#3}
  \end{subfigure}%
}
% #1--#4: dosya yolları (soldan sağa, üstten alta; boşsa yer tutucu);
% #5--#8: alt yazılar. \caption ve \label çağıran tarafta kalır -- bunu
% her zaman \begin{figure}[H] ... \end{figure} içine sarın (bkz. kullanım
% örnekleri deneyler.tex içinde \S{Gerçek Kayıtlar Üzerinde Girdi
% Konfigürasyonu}'nda).
\newcommand{\dortlugorsel}[8]{%
  \gorselhucre{#1}{#5}\hfill\gorselhucre{#2}{#6}\\[2mm]
  \gorselhucre{#3}{#7}\hfill\gorselhucre{#4}{#8}%
}

% ---- Bağlantılar, referanslar ------------------------------------------------
\usepackage{hyperref}
\hypersetup{
  colorlinks=true,
  linkcolor=anarenk,
  citecolor=yesilvurgu,
  urlcolor=vurgu,
  pdftitle={EdgeTAM'ın Jetson AGX Orin Üzerinde TensorRT ile Hızlandırılması},
  pdfauthor={Yiğit Kaya Bağcı},
}
% english: cleveref'in kendi dahili Türkçe ad tablosu polyglossia'nın
% Türkçe büyük/küçük harf kurallarıyla (noktasız ı / noktalı İ) çakışıp
% \begin{document}'te "Undefined control sequence" hatası veriyordu --
% bilinen bir cleveref+polyglossia-turkish etkileşim sorunu. english
% seçeneği cleveref'in dahili tablosunu devre dışı bırakır; görünen adlar
% zaten aşağıdaki \crefname tanımlarından geliyor.
\usepackage[capitalize, nameinlink, english]{cleveref}
\crefname{table}{Tablo}{Tablolar}
\Crefname{table}{Tablo}{Tablolar}
\crefname{figure}{Şekil}{Şekiller}
\Crefname{figure}{Şekil}{Şekiller}
\crefname{chapter}{Bölüm}{Bölümler}
\Crefname{chapter}{Bölüm}{Bölümler}
\crefname{section}{Bölüm}{Bölümler}
\Crefname{section}{Bölüm}{Bölümler}

% ---- Kaynakça -----------------------------------------------------------------
\usepackage[
  backend=biber,
  style=numeric,
  sorting=none,
  giveninits=true,
]{biblatex}
\addbibresource{kaynaklar.bib}

% ---- Matematik ------------------------------------------------------------
\usepackage{amsmath}
\usepackage{amssymb}

% \theauthor / \thedate: \title, \author, \date içeriğine \maketitle
% dışında da erişim (özel başlık sayfası için).
\usepackage{titling}

% ---- Üstbilgi/altbilgi ------------------------------------------------------
\usepackage[automark]{scrlayer-scrpage}
\pagestyle{scrheadings}
\clearpairofpagestyles
\ihead{\headmark}
\ohead{\thepage}
\ifoot{}
\ofoot{EdgeTAM / TensorRT -- Jetson AGX Orin}
